\documentclass[11pt,reqno]{amsart}

% --- Packages -----------------------------------------------------------------
\usepackage[T1]{fontenc}
\usepackage[utf8]{inputenc}
\usepackage[english]{babel}
\usepackage{amsmath,amssymb,amsthm,mathtools}
\usepackage{microtype}
\usepackage{booktabs}
\usepackage{enumitem}
\usepackage[dvipsnames]{xcolor}
\usepackage{hyperref}
\hypersetup{
  colorlinks=true,
  linkcolor=blue!50!black,
  citecolor=blue!50!black,
  urlcolor=blue!50!black,
}

% --- Theorem environments -----------------------------------------------------
\theoremstyle{plain}
\newtheorem{theorem}{Theorem}[section]
\newtheorem{lemma}[theorem]{Lemma}
\newtheorem{corollary}[theorem]{Corollary}
\newtheorem{proposition}[theorem]{Proposition}

\theoremstyle{definition}
\newtheorem{definition}[theorem]{Definition}

\theoremstyle{remark}
\newtheorem{remark}[theorem]{Remark}

% --- Macros -------------------------------------------------------------------
\newcommand{\Tplus}{T_{+}}
\newcommand{\Tminus}{T_{-}}
\newcommand{\Z}{\mathbb{Z}}
\newcommand{\N}{\mathbb{N}}
\newcommand{\Q}{\mathbb{Q}}
\newcommand{\Nodd}{\N_{\mathrm{odd}}}
\newcommand{\Obs}{\mathcal{O}}
\newcommand{\Atom}{\mathcal{A}}
\newcommand{\Fac}{\mathcal{F}}
\newcommand{\Xend}{X_{\mathrm{end}}}
\newcommand{\cend}{c_{\mathrm{end}}}
\newcommand{\dend}{d_{\mathrm{end}}}
\newcommand{\Vk}{V_{K}}
\newcommand{\Vi}{V_{I}}
\newcommand{\vk}{v_{K}}
\newcommand{\vi}{v_{I}}
\newcommand{\vtwo}{v_{2}}

% --- Document -----------------------------------------------------------------
\title[Obstruction residues for the $3x-1$ map]{Obstruction residues for the \texorpdfstring{$3x-1$}{3x-1} map:\\
  full factor complexity, positive density,\\ and a constructive lift}

\author{Jakob Stemberger}
\address{Independent researcher}
\email{[email]}
\date{\today}

\subjclass[2020]{11B83, 37P05, 68R15, 11A07}

\keywords{Collatz conjecture, $3x-1$ map, residue dynamics, factor complexity, subshift, parallel reduction}

\begin{document}

\begin{abstract}
We study a family $\Obs_L \subseteq \{1, \ldots, 2^L\}$ of residue classes
modulo $2^L$ arising naturally in the dynamics of the $3x-1$ map
$\Tminus(n) := (3n-1)/2^{\vtwo(3n-1)}$. The family is defined by an
algebraic criterion on a two-track parallel reduction.
We prove three structural results: a constructive lift between
modular levels, the existence of a positive asymptotic density
$c_W \in [3/64, 1]$, and full factor complexity $p_W(n) = 2^n$ of the
associated language of binary representations.
The central technical tool is a parity lemma in the residue system, which
forces the terminal data of the parallel reduction into a shift-$a$ normal
form with integer-valued shift index. A finer classification by shift
index yields a strict lift balance, an exact series representation of
$c_W$, and reduces the open question of $c_W$'s exact value to a single
quantitative bound on a sub-family of atomic obstructions. Results
transfer to the classical $3n+1$ map via the Lagarias intertwining.
\end{abstract}

\maketitle

% =============================================================================
\section{Introduction}
% =============================================================================

\subsection{The Collatz map and its \texorpdfstring{$3x-1$}{3x-1} cousin}

Let $\Nodd$ denote the set of positive odd integers, and let $\vtwo(m)$ denote
the $2$-adic valuation of an integer $m \ne 0$. The \emph{Collatz map}
$\Tplus$ takes an odd integer $n$ to the odd part of $3n+1$:
\[
\Tplus(n) := \frac{3n+1}{2^{\vtwo(3n+1)}}.
\]
The Collatz conjecture, attributed to Lothar Collatz and widely studied since
the 1930s, asserts that for every $n \in \Nodd$ the iteration
$n, \Tplus(n), \Tplus^2(n), \ldots$ eventually reaches the fixed point $1$.
Despite considerable computational evidence and a large body of partial
results --- most recently the breakthrough of Tao~\cite{Tao2022}, showing
that almost all trajectories attain almost bounded values --- the conjecture
itself remains open. We refer the reader to the annotated bibliographies
of Lagarias~\cite{Lagarias2010,Lagarias2010bib} for a comprehensive overview
of the literature.

The Collatz map has a natural cousin, the \emph{$3x-1$ map} $\Tminus$ acting
on the same domain:
\[
\Tminus(n) := \frac{3n-1}{2^{\vtwo(3n-1)}}.
\]
The two maps are structurally parallel: both perform the operation
``multiply by three, add a small constant, strip the factors of two from the
result''. They differ only in the sign of the additive constant. The
trajectories of $\Tminus$ are, however, \emph{not} believed to converge to a
single attractor: three nontrivial cycles are classically known --- the
fixed cycle through~$1$, the $\{5, 7\}$-cycle, and the seventeen-family cycle
through $\{17, 25, 37, 55, 65, 49, 73, \ldots\}$ --- and the question of
whether further cycles exist remains as open as the Collatz conjecture itself.
The two maps are connected through a classical identity going back to
Lagarias~\cite{Lagarias1985}: every $\Tminus$-trajectory starting from $m$
corresponds, via the substitution $r = 2m+1$, to a single $\Tplus$-step
followed by a $\Tplus$-trajectory. This \emph{Lagarias intertwining} means
that any structural result about the residue dynamics of $\Tminus$ has
implications for $\Tplus$, and conversely.

The present paper is concerned with a family of residue classes modulo
$2^L$ that arises naturally in the dynamics of $\Tminus$. We show that this
family --- which we shall call the \emph{obstruction residues} for reasons
made precise in Section~\ref{sec:setup} --- has three structural properties:

\begin{enumerate}[label=(\arabic*)]
\item a constructive lift between modular levels $L$ and $L+1$,
\item positive asymptotic density in $\{1, \ldots, 2^L\}$ as $L \to \infty$,
\item maximal combinatorial complexity of its associated language of bit
  strings.
\end{enumerate}

The third property is in some tension with the first two: a constructive
lift and positive density would, in many natural settings, force a certain
rigidity of the resulting word language (one might expect a sofic shift or a
primitive substitution shift). We prove that no such substitution-type
characterization is possible, identifying the obstruction residues as a
structurally richer family than the naive analogies would suggest.

\subsection{Parallel reductions and the \texorpdfstring{$X$}{X}-invariant}

The construction is built on a simple observation. For an odd integer $r$
with $v := \vtwo(r-1) \ge 1$, write $m := (r-1)/2^v$. We run
$\Tminus$ in parallel on $r$ and on $m$ modulo $2^L$ for a fixed level $L$,
keeping track at each step both of the current values and of their
\emph{affine relationship}: a pair $(c_j, d_j) \in \Q^2$ such that
$\Tminus^j(m) \equiv c_j \cdot \Tminus^j(r) + d_j$ modulo an
appropriately reduced power of two. The pair is updated by an explicit rule
depending on the $2$-adic valuations of the two trajectories. We denote the
terminal pair $(\cend, \dend)$ and define the scalar invariant
\[
\Xend(r, L) := 3 \dend(r, L) + \cend(r, L).
\]

For an odd integer $r$ with $\vtwo(r-1) \ge 1$, we
call $r$ an \emph{obstruction residue at level $L$} when $\Xend(r, L) = 1$,
and write $\Obs_L \subseteq \{1, \ldots, 2^L\}$ for the set of all such
residues. The precise definition is given in Section~\ref{sec:setup}
(Definition~\ref{def:obstruction}); the typical case is $v = 1$, i.e.\
$r \equiv 3 \pmod 4$.

The family $(\Obs_L)_{L \ge 6}$ is small at first --- $\Obs_6 = \{19, 27, 59\}$ ---
but grows rapidly: direct enumeration gives
\[
|\Obs_6| = 3, \quad |\Obs_7| = 8, \quad |\Obs_8| = 19, \quad \ldots, \quad |\Obs_{16}| = 7556,
\]
and the density ratio $|\Obs_L|/2^L$ rises monotonically. Three structural
questions arise:

\begin{enumerate}[label=(Q\arabic*)]
\item Does $r_0 \in \Obs_{L_0}$ guarantee that both $r_0$ and $r_0 + 2^{L_0}$
  lie in $\Obs_{L_0+1}$?
\item Does $|\Obs_L|/2^L$ converge to a positive limit, and can the limit
  be bounded from below?
\item What is the set of binary words appearing as factors of the bit
  representations of residues in $\bigcup_L \Obs_L$?
\end{enumerate}

The three questions are not independent: a positive answer to (Q1) gives
monotonicity, hence most of (Q2). A constructive form of (Q1), exhibiting
arbitrary bit patterns under lift, gives (Q3). All three are answered in
the affirmative by the main results.

\subsection{Main results}

\begin{theorem}[Constructive lift]
\label{thm:lift-intro}
Let $L_0 \ge 6$ and $r_0 \in \Obs_{L_0}$. Then both $r_+ := r_0$ and
$r_- := r_0 + 2^{L_0}$ are obstruction residues at level $L_0 + 1$.
\end{theorem}

The proof, given in Section~\ref{sec:lift}, proceeds by case analysis on a
finite invariant of the parallel reduction. (The hypothesis $L_0 \ge 6$ is
the smallest level for which $\Obs_{L_0}$ is non-empty; for $L_0 < 6$ the
statement is vacuous. The proof in Section~\ref{sec:lift} only uses
$L_0 > v$, which is automatic from $r_0 \in \Obs_{L_0}$ by
Definition~\ref{def:obstruction}.) The lift is explicit: for each $r_0$,
the terminal data of the reduction at $r_+$ or $r_-$ is determined by a
closed-form algebraic identity.

\begin{theorem}[Positive density]
\label{thm:density-intro}
The density ratio $|\Obs_L|/2^L$ is non-decreasing in $L$, and the limit
\[
c_W := \lim_{L \to \infty} \frac{|\Obs_L|}{2^L}
\]
exists in $[3/64, 1]$. For every $L_0 \ge 2$, the rigorous lower bound
$c_W \ge |\Obs_{L_0}|/2^{L_0}$ holds.
\end{theorem}

A direct computation at $L_0 = 16$, using only the algorithmic
$\Xend$-criterion of Section~\ref{sec:setup}, gives the lower bound
$c_W \ge 7556/65536 \approx 0.115$. The conservative bound $3/64$ obtained
from $L_0 = 6$ admits substantial improvement by raising $L_0$, but the
exact value of $c_W$ remains an open problem. We give in
Section~\ref{sec:Ga} a refined decomposition that exhibits $c_W$ as a
rigorously convergent series and isolates the remaining open question to a
single quantitative bound on a sub-family of obstructions.

For $r \in \Obs_L$, write $w_r \in \{0,1\}^L$ for the binary representation
of $r$ as an $L$-bit word, least significant bit first. The
\emph{factor language} of the obstruction family is
\[
\Fac(\Obs) := \{\, u \in \{0,1\}^* : u \text{ is a factor of some } w_r,\; r \in \Obs_L \text{ for some } L \,\},
\]
and the \emph{factor complexity function} is
$p_W(n) := |\Fac(\Obs) \cap \{0,1\}^n|$.

\begin{theorem}[Full factor complexity]
\label{thm:fc-intro}
For every $n \ge 1$,
\[
p_W(n) = 2^n.
\]
\end{theorem}

The proof, in Section~\ref{sec:fc}, is constructive: for each binary word
$u$, we exhibit an explicit residue $r(u)$ in some $\Obs_L$ having $u$ as a
factor. The construction is a direct iteration of the lift theorem.

A consequence of Theorem~\ref{thm:fc-intro} is a sharp structural rigidity
(Corollary~\ref{cor:no-substitution}): the obstruction family cannot be
described as a subshift of finite type, a sofic shift, or a primitive
substitution shift in the sense of Pansiot~\cite{Pansiot1984} or
Devyatov~\cite{Devyatov2015}.

\subsection{Strategy}

All three theorems rest on a single algebraic structure: the \emph{shift-$a$
normal form} of the terminal data. We prove in Section~\ref{sec:setup}
(Lemma~\ref{lem:parity}) that whenever the parallel reduction terminates
with $\Xend = 1$, the terminal pair has the form
\[
(\cend, \dend) = \left(4^a, \frac{1 - 4^a}{3}\right) \qquad \text{for some } a \in \Z.
\]
The integer $a$, which we call the \emph{shift index}, measures the
difference between the $2$-adic valuations of the two parallel trajectories.
Theorem~\ref{thm:lift-intro} is proved by case analysis on how the shift
index evolves under lift; Theorem~\ref{thm:fc-intro} is the iterative
consequence; Theorem~\ref{thm:density-intro} follows by monotonicity
together with a refined classification by shift index.

The shift-$a$ normal form is, to the best of our knowledge, new; in
particular, it differs from the Markov-operator setup of
Wirsching~\cite{Wirsching1998}, which works with an invariant
probability-density function rather than a terminal algebraic datum. The
remaining analysis proceeds by elementary case analysis and combinatorial
bookkeeping; the arguments do not depend on specialized Collatz machinery
beyond the Lagarias intertwining.

\subsection{Outline}

Section~\ref{sec:setup} introduces the parallel reduction, defines the
$X$-invariant rigorously, and proves the parity lemma. Section~\ref{sec:lift}
contains the proof of Theorem~\ref{thm:lift-intro}, split into the two
stop-type cases. Section~\ref{sec:density} derives
Theorem~\ref{thm:density-intro} as a direct corollary. Section~\ref{sec:fc}
proves Theorem~\ref{thm:fc-intro} and derives the structural rigidity
corollary. Section~\ref{sec:Ga} develops the classification by shift index,
proves the strict lift balance and series representation, and isolates the
remaining open problem. Section~\ref{sec:tplus} transfers the results to the
$3x+1$ side via the residue bijection. Section~\ref{sec:open} discusses open
questions. Appendix~\ref{sec:appendix} summarizes computational
verifications.

% =============================================================================
\section{The parallel reduction and the \texorpdfstring{$X$}{X}-invariant}
\label{sec:setup}
% =============================================================================

\subsection{The basic step}

Recall that for an odd integer $n$, the map $\Tminus$ strips the $2$-adic
part from $3n - 1$:
\[
\Tminus(n) = \frac{3n - 1}{2^{\vtwo(3n - 1)}}.
\]
We need a version of this step that keeps track of the modulus. Fix a level
$L \ge 1$, and consider a pair $(\alpha, \beta)$ with $\alpha$ a positive odd
integer and $\beta$ a power of two, interpreted as a residue $\alpha$ modulo
$\beta$. A single $\Tminus$-step on such a pair produces
\[
(\alpha', \beta') = \left(\frac{3\alpha - 1}{2^{v_\alpha}},\; \frac{3\beta}{2^{v_\alpha}}\right), \qquad v_\alpha := \vtwo(3\alpha - 1),
\]
provided $v_\alpha < \vtwo(\beta)$. If $v_\alpha \ge \vtwo(\beta)$, the
modulus has been exhausted and the step is undefined; we say the reduction
has \emph{terminated}.

Since $3$ is odd, $\vtwo(3\beta) = \vtwo(\beta)$, so
$\vtwo(\beta') = \vtwo(\beta) - v_\alpha$: the modulus strictly decreases at
each step, and the reduction always terminates after finitely many steps.

\subsection{The two-track setup}

Fix $r$ a positive odd integer, set $v := \vtwo(r-1) \ge 1$, and write
$m := (r - 1)/2^v$; then $m$ is odd. (We use $v$ as a parameter throughout;
the special case $v = 1$ — equivalently $r \equiv 3 \pmod 4$ — is the
typical reader picture and is presented first wherever a concrete
calculation helps.)

Fix a level $L > v$. We run two $\Tminus$-reductions in parallel:
\begin{itemize}
\item the \emph{K-track}, starting from $(a_K^{(0)}, b_K^{(0)}) := (r, 2^L)$,
\item the \emph{I-track}, starting from $(a_I^{(0)}, b_I^{(0)}) := (m, 2^{L-v})$.
\end{itemize}
The initial relation $m = (r-1)/2^v$ writes as $a_I^{(0)} = 2^{-v} a_K^{(0)} - 2^{-v}$.
At each step $j \ge 0$, we update both tracks according to the basic step,
denoting $\vk^{(j)} := \vtwo(3 a_K^{(j)} - 1)$ and $\vi^{(j)} := \vtwo(3 a_I^{(j)} - 1)$.
The reduction proceeds as long as both tracks can step; it terminates as
soon as either condition fails. We write $J = J(r, L)$ for the index of
termination.

\subsection{The affine relation and the \texorpdfstring{$X$}{X}-invariant}

\begin{lemma}\label{lem:affine}
For each $j$ with $0 \le j \le J$, there exist $c_j, d_j \in \Q$ such that
\[
a_I^{(j)} = c_j \, a_K^{(j)} + d_j \qquad \text{in } \Q,
\]
with $c_0 = 2^{-v}$, $d_0 = -2^{-v}$, and the update rule
\[
c_{j+1} = c_j \cdot 2^{\vk^{(j)} - \vi^{(j)}}, \qquad
d_{j+1} = \frac{3 d_j + c_j - 1}{2^{\vi^{(j)}}}.
\]
\end{lemma}

\begin{proof}
The base case is immediate. Assume the affine relation holds at index $j$.
Writing $A_K^{(j)} := 3 a_K^{(j)} - 1$, $A_I^{(j)} := 3 a_I^{(j)} - 1$, and
$X_j := 3 d_j + c_j$,
\[
A_I^{(j)} = 3(c_j a_K^{(j)} + d_j) - 1 = c_j (3 a_K^{(j)} - 1) + (3 d_j + c_j - 1) = c_j A_K^{(j)} + (X_j - 1).
\]
Dividing by $2^{\vi^{(j)}}$ and using $A_K^{(j)} = 2^{\vk^{(j)}} a_K^{(j+1)}$:
\[
a_I^{(j+1)} = c_j \cdot 2^{\vk^{(j)} - \vi^{(j)}} a_K^{(j+1)} + \frac{X_j - 1}{2^{\vi^{(j)}}},
\]
which is the desired affine relation at index $j+1$. \qedhere
\end{proof}

We define the \emph{$X$-invariant} of the parallel reduction at $(r, L)$ as
\[
\Xend(r, L) := X_J = 3 d_J + c_J.
\]
This is not invariant under the update step; the name reflects its role as
the algebraic certificate of obstruction status (Definition~\ref{def:obstruction}).

\subsection{Definition of obstruction residues}

\begin{definition}[Obstruction residue]
\label{def:obstruction}
Let $r$ be a positive odd integer and write $v := \vtwo(r-1) \ge 1$. We
call $r$ an \emph{obstruction residue at level $L > v$} if $\Xend(r, L) = 1$.
We write
\[
\Obs_L := \{\, r \in \{1, \ldots, 2^L\} : r \text{ odd, and an obstruction residue at level } L \,\}.
\]
\end{definition}

The definition is algorithmic: given $r$ and $L$, one runs the parallel
reduction, observes the terminal pair $(c_J, d_J)$, and checks
$3 d_J + c_J = 1$. The computation terminates in at most $L$ steps.

The connection between this algebraic definition and the dynamics of
$\Tminus$ trajectories is, briefly, that $\Xend = 1$ implies a
synchronization of stopping times: starting from $r$ and from any
$n \equiv r \pmod{2^L}$, the stopping-time difference between the
$\Tminus$ and $\Tminus^{(3)}$ iterates --- where $\Tminus^{(3)}$ denotes the
$3$-applications-of-$\Tminus$ variant --- remains constant. We do not
develop this here; it is discussed in Section~\ref{sec:tplus}.

\subsection{The parity lemma}

The central algebraic fact of this paper is the following.

\begin{lemma}[Parity lemma]
\label{lem:parity}
Suppose the parallel reduction at $(r, L)$ terminates with $\Xend(r, L) = 1$.
Then
\[
\delta_J := \Vk^{(J)} - \Vi^{(J)} - v
\]
is an \emph{even} integer (possibly negative), where $v := \vtwo(r-1)$,
$\Vk^{(j)} := \sum_{i < j} \vk^{(i)}$ and analogously for $\Vi$.
\end{lemma}

\begin{proof}
Iterating the update rule from Lemma~\ref{lem:affine} with $c_0 = 2^{-v}$:
\[
c_j = 2^{-v} \cdot \prod_{i < j} 2^{\vk^{(i)} - \vi^{(i)}} = 2^{\Vk^{(j)} - \Vi^{(j)} - v} = 2^{\delta_j}.
\]
From $X_J = 3 d_J + c_J = 1$, $d_J = (1 - c_J)/3 = (1 - 2^{\delta_J})/3$.
The affine relation at index $J$ becomes
\[
a_I^{(J)} = 2^{\delta_J} a_K^{(J)} + \frac{1 - 2^{\delta_J}}{3} \qquad \text{in } \Q.
\]
Multiplying by $3$:
\begin{equation}\label{eq:parity-master}
3 a_I^{(J)} = 3 \cdot 2^{\delta_J} a_K^{(J)} + 1 - 2^{\delta_J}.
\end{equation}
We split by sign of $\delta_J$. Write $\delta := \delta_J$ throughout.

\emph{Case $\delta \ge 0$.} All terms in~\eqref{eq:parity-master} are integers.
Reducing modulo~$3$: $3 a_I^{(J)} \equiv 0$ and $3 \cdot 2^\delta a_K^{(J)} \equiv 0$,
so $1 - 2^\delta \equiv 0 \pmod 3$, i.e.\ $2^\delta \equiv 1 \pmod 3$. Since
$2 \equiv -1 \pmod 3$, this forces $\delta$ even.

\emph{Case $\delta < 0$.} Write $\delta = -e$, $e > 0$. Multiplying~\eqref{eq:parity-master}
by $2^e$:
\[
3 \cdot 2^e \, a_I^{(J)} = 3 \, a_K^{(J)} + 2^e - 1.
\]
All terms are integers; reducing modulo~$3$ yields $2^e - 1 \equiv 0 \pmod 3$,
hence $e$ even, hence $\delta$ even.

In both cases $\delta$ is even. \qedhere
\end{proof}

\subsection{The shift-\texorpdfstring{$a$}{a} normal form}

Writing $\delta_J = 2a$ with $a \in \Z$, the parity lemma yields
\[
c_J = 4^a, \qquad d_J = \frac{1 - 4^a}{3}.
\]
We call this the \emph{shift-$a$ normal form} and $a$ the \emph{shift index}
of $r$ at level $L$, denoted $a(r, L)$. The obstruction family partitions:
\[
\Obs_L = \bigsqcup_{a \in \Z} \Obs_L^{G_a}, \qquad \Obs_L^{G_a} := \{\, r \in \Obs_L : a(r, L) = a \,\}.
\]
For $a = 0$ we have $(c_J, d_J) = (1, 0)$, meaning $a_I^{(J)} = a_K^{(J)}$.
The shift index is bounded: $|2a + v| \le \max(\Vk^{(J)}, \Vi^{(J)}) \le L$,
so $|a| \le L/2$.

\subsection{Simultaneous stop and stop types}

\begin{corollary}[Simultaneous stop]
\label{cor:sync}
If $\Xend(r, L) = 1$, then the K-track and I-track satisfy the termination
condition $v_\alpha \ge \vtwo(\beta)$ at the same index $J$.
\end{corollary}

\begin{proof}
From the shift-$a$ normal form, $a_I^{(J)} = 4^a a_K^{(J)} + (1-4^a)/3$, so
$A_I^{(J)} = 4^a A_K^{(J)}$ and hence $\vi^{(J)} = \vk^{(J)} + 2a$.
Combining with $\vtwo(b_I^{(J)}) = (L-v) - \Vi^{(J)}$ and
$\vtwo(b_K^{(J)}) = L - \Vk^{(J)}$, and using $\Vk - \Vi = 2a + v$:
\begin{align*}
\vi^{(J)} \ge \vtwo(b_I^{(J)})
  &\iff \vk^{(J)} + 2a \ge L - v - \Vi^{(J)} \\
  &\iff \vk^{(J)} \ge L - \Vk^{(J)} \\
  &\iff \vk^{(J)} \ge \vtwo(b_K^{(J)}). \qedhere
\end{align*}
\end{proof}

The \emph{stop type} at $(r, L)$ is \emph{EE} (equality) if
$\vk^{(J)} = \vtwo(b_K^{(J)})$ (and hence $\vi^{(J)} = \vtwo(b_I^{(J)})$),
and \emph{SS} (strict) if both inequalities are strict.

\subsection{The parameter \texorpdfstring{$v$}{v} and the typical case \texorpdfstring{$v = 1$}{v = 1}}

The constructions above are stated for general $v := \vtwo(r-1) \ge 1$. The
typical case is $v = 1$ (equivalently $r \equiv 3 \pmod 4$); for $r \equiv 1
\pmod 4$ the value of $v$ is $\ge 2$. In every formula above, setting $v = 1$
recovers the original initial pair $(c_0, d_0) = (\tfrac{1}{2}, -\tfrac{1}{2})$
and the parity relation $\Vk^{(J)} - \Vi^{(J)} \equiv 1 \pmod 2$. The
shift-$a$ normal form $(c_J, d_J) = (4^a, (1-4^a)/3)$ is independent of $v$.
The subsequent sections work with the general case throughout; explicit
calculations use $v = 1$ wherever this aids the reader.

The $v \ge 2$ residues contribute a small but non-negligible share of
$\Obs_L$: direct enumeration gives, for $L = 16$, $|\Obs_L| = 7556$ in
total, of which $6130$ have $v = 1$ and the remaining $1426$ (roughly
$19\%$) have $v \ge 2$. The numbers in the appendix table count all $v$.

% =============================================================================
\section{The Lift Theorem}
\label{sec:lift}
% =============================================================================

We now prove Theorem~\ref{thm:lift-intro}: for every $r_0 \in \Obs_{L_0}$,
both lifts $r_+ := r_0$ and $r_- := r_0 + 2^{L_0}$ are obstructions at level
$L_0 + 1$. Throughout, $a \in \Z$ is the shift index of $r_0$ at level $L_0$,
and we use the synchronization $\vi^{(J)} = \vk^{(J)} + 2a$
(Corollary~\ref{cor:sync}) freely.

\subsection{The \texorpdfstring{$r_+$}{r+} lift}

\begin{theorem}\label{thm:lift-plus}
Let $r_0 \in \Obs_{L_0}^{G_a}$. Then $r_+ := r_0 \in \Obs_{L_0+1}$.
\end{theorem}

\begin{proof}
Set $v := \vtwo(r_0 - 1) \ge 1$.
The starting pairs at $(r_+, L_0+1)$ are $(r_0, 2^{L_0+1})$ and
$(m_0, 2^{L_0+1-v})$ where $m_0 = (r_0 - 1)/2^v$. These differ from the
corresponding pairs at $(r_0, L_0)$ only in the modulus, which is doubled.
Since each basic step depends only on $a_K^{(j)}, a_I^{(j)}$ and not on
the modulus (as long as the latter does not constrain the step), the
parallel reduction at $(r_+, L_0+1)$ produces identical
$\vk^{(j)}, \vi^{(j)}, c_j, d_j$ for $j \le J$ to those at $(r_0, L_0)$.

\emph{Stop type EE at $L_0$.} Here $\vk^{(J)} = L_0 - \Vk^{(J)}$ exactly. At
level $L_0 + 1$, the K-track modulus at index $J$ has $2$-adic valuation
$L_0 + 1 - \Vk^{(J)} = \vk^{(J)} + 1 > \vk^{(J)}$, so the step is permitted.
Taking one more step, with $c_J = 4^a$, $d_J = (1 - 4^a)/3$:
\[
c_{J+1} = 4^a \cdot 2^{\vk^{(J)} - \vi^{(J)}} = 4^a \cdot 2^{-2a} = 1,
\]
\[
d_{J+1} = \frac{3(1 - 4^a)/3 + 4^a - 1}{2^{\vi^{(J)}}} = 0.
\]
After this step, $\Vk^{(J+1)} = L_0$, so the K-track modulus has valuation
$1$, and the next step is impossible (it would require $\vk^{(J+1)} < 1$,
impossible for odd $a_K^{(J+1)}$). Termination occurs at index $J+1$ with
$(c_{J+1}, d_{J+1}) = (1, 0)$, hence $\Xend = 1$ and
$r_+ \in \Obs_{L_0+1}^{G_0}$.

\emph{Stop type SS at $L_0$.} Here $\vk^{(J)} > L_0 - \Vk^{(J)}$, i.e.\
$\vk^{(J)} \ge L_0 + 1 - \Vk^{(J)}$ (integer), which is exactly the
termination condition at level $L_0 + 1$. Termination at $(r_+, L_0+1)$
occurs at index $J$ with terminal data unchanged, so
$r_+ \in \Obs_{L_0+1}^{G_a}$. \qedhere
\end{proof}

\subsection{The \texorpdfstring{$r_-$}{r-} lift}

\begin{theorem}\label{thm:lift-minus}
Let $r_0 \in \Obs_{L_0}^{G_a}$. Then $r_- := r_0 + 2^{L_0} \in \Obs_{L_0+1}$.
\end{theorem}

\begin{proof}
Write $a_{K,+}^{(j)}, a_{I,+}^{(j)}$ for the values at $(r_+, L_0+1)$
(from Theorem~\ref{thm:lift-plus}) and $a_{K,-}^{(j)}, a_{I,-}^{(j)}$ for
those at $(r_-, L_0+1)$. The discrepancies
\[
\Delta_K^{(j)} := a_{K,-}^{(j)} - a_{K,+}^{(j)}, \qquad \Delta_I^{(j)} := a_{I,-}^{(j)} - a_{I,+}^{(j)}
\]
satisfy $\Delta_K^{(0)} = 2^{L_0}$, $\Delta_I^{(0)} = 2^{L_0 - v}$ (the
I-track shift $\Delta_I^{(0)} = (r_- - 1)/2^v - m_0 = 2^{L_0}/2^v$ uses
$\vtwo(r_- - 1) = \min(\vtwo(r_0 - 1), L_0) = v$, valid for $L_0 > v$).
We claim that as long as $\Vk^{(j)} < L_0$ and $\Vi^{(j)} < L_0 - v$, the two
reductions take identical local steps,
\begin{equation}\label{eq:lift-minus-vmatch}
v_2\!\left(3 a_{K,-}^{(j)} - 1\right) = v_2\!\left(3 a_{K,+}^{(j)} - 1\right) = \vk^{(j)}
\end{equation}
(and analogously for the I-track), and the discrepancies have the closed
form
\begin{equation}\label{eq:lift-minus-delta}
\Delta_K^{(j)} = 2^{L_0 - \Vk^{(j)}} \cdot 3^j, \qquad \Delta_I^{(j)} = 2^{L_0 - v - \Vi^{(j)}} \cdot 3^j.
\end{equation}

We prove \eqref{eq:lift-minus-vmatch} and \eqref{eq:lift-minus-delta} jointly
by induction on $j$ in the range $0 \le j < J$, where $J$ is the termination
index of the reduction at $(r_0, L_0)$. The base case $j = 0$ is immediate.
For the inductive step, assume \eqref{eq:lift-minus-delta} at step $j$ and
$j < J$. Since the K-track at $(r_0, L_0)$ has not yet terminated,
$\vk^{(j)} < L_0 - \Vk^{(j)} = \vtwo(\Delta_K^{(j)})$. The decomposition
$3 a_{K,-}^{(j)} - 1 = (3 a_{K,+}^{(j)} - 1) + 3 \Delta_K^{(j)}$ then has its
two summands at strictly different valuations $\vk^{(j)}$ and
$L_0 - \Vk^{(j)}$, so by the ultrametric identity
$\vtwo(x + y) = \min(\vtwo x, \vtwo y)$ when $\vtwo x \ne \vtwo y$, the sum
has valuation exactly $\vk^{(j)}$. This proves \eqref{eq:lift-minus-vmatch}
at step $j$. The basic step then yields
$\Delta_K^{(j+1)} = (3 \Delta_K^{(j)})/2^{\vk^{(j)}} = 2^{L_0 - \Vk^{(j+1)}} \cdot 3^{j+1}$;
the same calculation on the I-track completes
\eqref{eq:lift-minus-delta} at step $j+1$.

\emph{Stop type SS at $L_0$.} At index $J$ of $(r_+, L_0+1)$,
$\vk^{(J)} > L_0 - \Vk^{(J)}$, hence $\vtwo(\Delta_K^{(J)}) = L_0 - \Vk^{(J)} < \vk^{(J)}$.
The relation
$3 a_{K,-}^{(J)} - 1 = (3 a_{K,+}^{(J)} - 1) + 3 \Delta_K^{(J)}$ now has the
two summands at distinct valuations $\vk^{(J)}$ and $L_0 - \Vk^{(J)}$, so by
the ultrametric inequality the sum has valuation
$\vk^{(J),-} := L_0 - \Vk^{(J)} < L_0 + 1 - \Vk^{(J)}$. The K-step at
$(r_-, L_0+1)$ at index $J$ is therefore permitted. The same calculation
on the I-track --- using $\vtwo(\Delta_I^{(J)}) = L_0 - v - \Vi^{(J)} < \vi^{(J)}$,
which holds by SS on the I-track via the synchronisation
$\vi^{(J)} = \vk^{(J)} + 2a$ of Corollary~\ref{cor:sync} --- yields
$\vi^{(J),-} := L_0 - v - \Vi^{(J)}$, also permitted, and one more step
follows. Using $\Vk - \Vi = 2a + v$, the new exponent in the update rule
is $\vk^{(J),-} - \vi^{(J),-} = -2a$, so
\[
c_{J+1} = c_J \cdot 2^{-2a} = 4^a \cdot 4^{-a} = 1,
\qquad
d_{J+1} = \frac{3 \cdot (1 - 4^a)/3 + 4^a - 1}{2^{\vi^{(J),-}}} = 0.
\]
After this step $\Vk^{(J+1)} = L_0$ and $\Vi^{(J+1)} = L_0 - v$, so both
moduli have valuation $1$ and the next step on either track is impossible
for odd $a_K, a_I$. Termination occurs at index $J + 1$ with
$(c_{J+1}, d_{J+1}) = (1, 0)$, so $r_- \in \Obs_{L_0+1}^{G_0}$.

\emph{Stop type EE at $L_0$.} Here
$\vk^{(J)} = L_0 - \Vk^{(J)} = \vtwo(\Delta_K^{(J)})$, so the two summands
of $3 a_{K,-}^{(J)} - 1 = (3 a_{K,+}^{(J)} - 1) + 3 \Delta_K^{(J)}$ have
\emph{equal} valuation $\vk^{(J)}$. The ultrametric inequality then gives
only $\vtwo(A_{K,-}^{(J)}) \ge \vk^{(J)} + 1$; that is,
$\vtwo(A_{K,-}^{(J)}) \ge L_0 + 1 - \Vk^{(J)}$, which is the K-track stop
condition at level $L_0+1$. The same argument on the I-track gives
$\vtwo(A_{I,-}^{(J)}) \ge L_0 + 1 - v - \Vi^{(J)}$, the I-track stop
condition at level $L_0+1$. Termination therefore occurs at index $J$ with
terminal data unchanged, so $r_- \in \Obs_{L_0+1}^{G_a}$. \qedhere
\end{proof}

\subsection{The symmetry table}

The proofs of Theorems~\ref{thm:lift-plus} and~\ref{thm:lift-minus}
establish the following symmetry, central for the recurrence developed in
Section~\ref{sec:Ga}:

\begin{center}
\begin{tabular}{c|cc}
\toprule
Stop type of $r_0$ at $L_0$ & shift index of $r_+$ & shift index of $r_-$ \\
\midrule
EE & $0$ & $a$ \\
SS & $a$ & $0$ \\
\bottomrule
\end{tabular}
\end{center}
(Both lifts are at level $L_0 + 1$.)

So lifting an obstruction with shift index $a \ne 0$ always produces one
shift-$a$ obstruction and one shift-zero obstruction at the next level;
lifting a shift-zero obstruction produces two shift-zero obstructions.

% =============================================================================
\section{Positive density}
\label{sec:density}
% =============================================================================

\begin{proof}[Proof of Theorem~\ref{thm:density-intro}]
By Theorem~\ref{thm:lift-intro}, every $r_0 \in \Obs_{L_0}$ contributes two
distinct elements $r_0$ and $r_0 + 2^{L_0}$ to $\Obs_{L_0+1}$. Hence
$|\Obs_{L_0+1}| \ge 2 |\Obs_{L_0}|$, equivalently
$|\Obs_{L_0+1}|/2^{L_0+1} \ge |\Obs_{L_0}|/2^{L_0}$: the density ratio is
non-decreasing. Since $\Obs_L \subseteq \{1, \ldots, 2^L\}$ the ratio is
bounded above by $1$. By monotone convergence,
$c_W := \lim_{L \to \infty} |\Obs_L|/2^L$ exists in $[\rho_{L_0}, 1]$ for any
$L_0 \ge 2$. Specialising to $L_0 = 6$ and using the direct enumeration
$\Obs_6 = \{19, 27, 59\}$ (Section~\ref{sec:setup}; verified computationally
in Appendix~\ref{sec:appendix}), $c_W \ge 3/64$. \qedhere
\end{proof}

The conservative bound $\rho_6 = 3/64$ uses only the three explicit
obstructions $\Obs_6 = \{19, 27, 59\}$. By raising $L_0$ and computing
$|\Obs_{L_0}|$ directly, tighter rigorous bounds for $c_W$ follow:

\begin{center}
\begin{tabular}{rrrr}
\toprule
$L_0$ & $|\Obs_{L_0}|$ & $\rho_{L_0}$ & rigorous lower bound for $c_W$ \\
\midrule
$6$  & $3$    & $\approx 0.047$ & $3/64$ \\
$10$ & $91$   & $\approx 0.089$ & $91/1024$ \\
$14$ & $1776$ & $\approx 0.108$ & $1776/16384$ \\
$16$ & $7556$ & $\approx 0.115$ & $7556/65536$ \\
\bottomrule
\end{tabular}
\end{center}

These are not asymptotic estimates: they are finite-level computations
producing rigorous bounds at the cost of $O(L_0 \cdot 2^{L_0})$ operations
to enumerate obstructions modulo $2^{L_0}$. The exact value of $c_W$ is
not pinned down by this approach; we return to it in
Section~\ref{sec:Ga}.

% =============================================================================
\section{Full factor complexity}
\label{sec:fc}
% =============================================================================

\subsection{Setup}

To each $r \in \Obs_L$ we associate its binary representation as an
$L$-bit word, least significant bit first:
\[
w_r \in \{0,1\}^L, \qquad w_r[k] := \lfloor r / 2^k \rfloor \bmod 2, \quad k = 0, \ldots, L-1.
\]
The \emph{language of obstruction words} at level $L$ is
$\Fac_L := \{w_r : r \in \Obs_L\}$, and the \emph{factor language} of the
obstruction family is
\[
\Fac(\Obs) := \{\, u \in \{0,1\}^* : u \text{ is a factor of } w_r \text{ for some } r \in \Obs_L,\, L \ge 1 \,\}.
\]
The \emph{factor complexity function} is
$p_W(n) := |\Fac(\Obs) \cap \{0,1\}^n|$. Trivially $p_W(n) \le 2^n$.

\subsection{The main theorem}

\begin{theorem}\label{thm:fc}
For every $n \ge 1$, $p_W(n) = 2^n$.
\end{theorem}

\begin{proof}
Fix $n \ge 1$ and $u = (u_0, u_1, \ldots, u_{n-1}) \in \{0,1\}^n$. Choose
any anchor $r_0 \in \Obs_6 = \{19, 27, 59\}$ (verified by direct enumeration;
see Appendix~\ref{sec:appendix}). Define inductively
\[
r_k := r_{k-1} + u_{k-1} \cdot 2^{5 + k}, \qquad k = 1, 2, \ldots, n.
\]
We claim $r_k \in \Obs_{6+k}$ for each $k$. The base $k = 0$ is by choice of
$r_0$. For the inductive step from $k-1$ to $k$: if $u_{k-1} = 0$ then
$r_k = r_{k-1}$, so $r_k$ is the $r_+$ lift and lies in $\Obs_{6+k}$ by
Theorem~\ref{thm:lift-plus}; if $u_{k-1} = 1$ then $r_k = r_{k-1} + 2^{5+k}$
is the $r_-$ lift and lies in $\Obs_{6+k}$ by Theorem~\ref{thm:lift-minus}.

By construction, $w_{r_n}[5+k] = u_{k-1}$ for $k = 1, \ldots, n$, so $u$
appears as a factor of $w_{r_n}$ at positions $6, \ldots, 5+n$. Thus
$u \in \Fac(\Obs)$ for every $u \in \{0,1\}^n$, hence $p_W(n) \ge 2^n$.
Combined with the trivial upper bound, $p_W(n) = 2^n$. \qedhere
\end{proof}

\subsection{Entropy and substitution-shift rigidity}

The \emph{topological entropy} of a language $\mathcal{L} \subseteq \{0,1\}^*$
is $h(\mathcal{L}) := \limsup_{n \to \infty} (1/n) \log p_{\mathcal{L}}(n)$;
see Lind--Marcus~\cite{LindMarcus1995} for general background on entropy of
subshifts and subshifts of finite type, and Fogg~\cite{Fogg2002} or
Allouche--Shallit~\cite{AlloucheShallit2003} for factor complexity of
substitution and automatic sequences.

\begin{corollary}\label{cor:entropy}
The topological entropy of $\Fac(\Obs)$ is $\log 2$.
\end{corollary}

This is the maximal value, matching that of the full Bernoulli shift
$\{0,1\}^\Z$.

\begin{corollary}
\label{cor:no-substitution}
The obstruction family $(\Obs_L)_{L \ge 6}$, viewed as a graded family of
finite languages over $\{0,1\}$, is
\begin{enumerate}[label=(\roman*)]
\item not a non-trivial subshift of finite type;
\item not a non-trivial sofic shift in the classical sense;
\item not a primitive substitution shift in the sense of Pansiot~\cite{Pansiot1984}
  or Devyatov~\cite{Devyatov2015}.
\end{enumerate}
\end{corollary}

\begin{proof}
A non-trivial SFT defined by a non-empty finite list of forbidden subwords
has entropy strictly less than $\log 2$ (Lind--Marcus~\cite{LindMarcus1995}),
contradicting Corollary~\ref{cor:entropy}. A non-trivial sofic shift is the image of an
SFT under a one-block factor map and inherits the entropy constraint.
Primitive substitution shifts have factor complexity in one of the polynomial
classes $\Theta(1)$, $\Theta(n)$, $\Theta(n \log \log n)$, $\Theta(n \log n)$,
$\Theta(n^2)$ (Pansiot~\cite{Pansiot1984}; refined by
Devyatov~\cite{Devyatov2015}), giving entropy zero --- but our
$p_W(n) = 2^n$ has positive entropy. \qedhere
\end{proof}

% =============================================================================
\section{The \texorpdfstring{$G_a$}{Ga} classification}
\label{sec:Ga}
% =============================================================================

\subsection{Atomic obstructions}

An obstruction $r \in \Obs_L$ is \emph{atomic at level $L$} if
$r \bmod 2^{L-1} \notin \Obs_{L-1}$. We write $\Atom_L \subseteq \Obs_L$ and
$\Atom_L^{G_a} := \Atom_L \cap \Obs_L^{G_a}$.

\begin{lemma}[Atom decomposition]
\label{lem:atom-decomp}
For every $L \ge 6$,
\[
|\Obs_L| = \sum_{L_0 = 6}^{L} |\Atom_{L_0}| \cdot 2^{L - L_0}.
\]
\end{lemma}

\begin{proof}
We show that every $r \in \Obs_L$ has a unique smallest $L_0 \ge 6$ with
$r \bmod 2^{L_0} \in \Atom_{L_0}$. \emph{Existence:} starting from $r$ at
level $L$, consider the descending sequence $r_k := r \bmod 2^{L-k}$ for
$k = 0, 1, \ldots, L - 6$. By the lift theorem (Theorem~\ref{thm:lift-intro})
each $r_{k+1} \in \Obs_{L-k-1}$ implies $r_k \in \Obs_{L-k}$, but the converse
need not hold; let $L_0$ be the smallest level $\ge 6$ for which $r \bmod 2^{L_0} \in \Obs_{L_0}$.
Then $r \bmod 2^{L_0 - 1} \notin \Obs_{L_0 - 1}$ (either because $L_0 = 6$ and
$\Obs_5 = \emptyset$, or by minimality), so $r \bmod 2^{L_0} \in \Atom_{L_0}$.
\emph{Uniqueness:} the level $L_0$ is the smallest such by definition, hence
unique. Given the atomic anchor $r_0 := r \bmod 2^{L_0}$, the residue $r$ is
determined by $r_0$ together with the bits $\{u_{L_0}, \ldots, u_{L-1}\}$
encoding the lift choices (Theorems~\ref{thm:lift-plus},~\ref{thm:lift-minus}).
Each atomic anchor at level $L_0$ thus generates $2^{L - L_0}$ obstructions
at level $L$. \qedhere
\end{proof}

\subsection{No shift-zero atoms}

\begin{theorem}\label{thm:no-G0-atom}
For every $L \ge 7$, $|\Atom_L^{G_0}| = 0$.
\end{theorem}

\begin{proof}
Fix $L \ge 7$ and suppose $r \in \Obs_L^{G_0}$, so $\Xend(r, L) = 1$ with
shift index $a = 0$, i.e.\ $(c_J, d_J) = (1, 0)$. Write $r' := r \bmod 2^{L-1}$.
The parallel reduction at $(r', L-1)$ uses the same starting values
$a_K, a_I$ as at $(r, L)$, but with the modulus reduced by one power of two.
The first $j$ steps of the two reductions coincide as long as the modulus
does not constrain a step.

Termination of the reduction at $(r, L)$ at index $J$ means at least one
track satisfies its stop condition there; by Corollary~\ref{cor:sync}, both
do. Since $a = 0$, the synchronisation $\vi^{(J)} = \vk^{(J)} + 2a = \vk^{(J)}$
and $\Vk^{(J)} = \Vi^{(J)} + 1$ hold, so the K-track stop condition
$\vk^{(J)} \ge L - \Vk^{(J)}$ and the I-track stop condition
$\vi^{(J)} \ge (L-1) - \Vi^{(J)}$ are equivalent: both reduce to
$\Vk^{(J)} + \vk^{(J)} \ge L$. We split by whether the level-$(L-1)$ stop
condition $\Vk^{(j)} + \vk^{(j)} \ge L - 1$ first fires at index $J$ or
already at index $J - 1$. These are the only possibilities: if it fired at
some $J' \le J - 2$, then $\Vk^{(J')} + \vk^{(J')} \ge L - 1$; the
$(r, L)$ reduction takes its $J'$-th step (the level-$L$ stop has not yet
fired since $J' < J$), giving $\Vk^{(J'+1)} = L - 1$, so the next K-step
requires only $\vk^{(J'+1)} \ge 1$, which is automatic for odd
$a_K^{(J'+1)}$. Hence the $(r, L)$ reduction stops at $J' + 1 \le J - 1$,
contradicting the definition of $J$.

\emph{Case A: $\Vk^{(J)} + \vk^{(J)} \ge L$ and the reduction at
$(r', L-1)$ also terminates at index $J$.} Then the terminal data
coincide: $(c_J, d_J) = (1, 0)$, so $r' \in \Obs_{L-1}^{G_0}$ and $r$ is
not atomic.

\emph{Case B: $\Vk^{(J-1)} + \vk^{(J-1)} \ge L - 1$ but $< L$.} The parallel
reduction at $(r', L-1)$ terminates one step earlier, at index $J - 1$,
with terminal data $(c_{J-1}, d_{J-1})$. From the update rule
(Lemma~\ref{lem:affine}) at index $J - 1$,
\[
c_J = c_{J-1} \cdot 2^{\vk^{(J-1)} - \vi^{(J-1)}}, \qquad d_J = \frac{3 d_{J-1} + c_{J-1} - 1}{2^{\vi^{(J-1)}}}.
\]
Setting $c_J = 1$, $d_J = 0$ and reading off the numerator of $d_J$:
$3 d_{J-1} + c_{J-1} - 1 = 0$, i.e.\ $X_{J-1} = 1$. By the parity lemma
applied at the terminal index $J - 1$ of the reduction at $(r', L-1)$,
the terminal pair is in shift-$a'$ normal form for some $a' \in \Z$, so
$r' \in \Obs_{L-1}^{G_{a'}}$.

In either case, $r' \in \Obs_{L-1}$, so $r$ is not atomic. \qedhere
\end{proof}

\subsection{Strict lift balance and series representation}

Write $g_L := |\Obs_L^{G_0}|$ and $h_L := |\Obs_L^{G_{\ne 0}}|$. The lift
symmetry table from Section~\ref{sec:lift} gives, per parent at $L - 1$:

\begin{center}
\begin{tabular}{c|cc}
\toprule
Parent at $L-1$ & Contribution to $g_L$ & Contribution to $h_L$ \\
\midrule
$G_0$ & $2$ & $0$ \\
$G_a$, $a \ne 0$ & $1$ & $1$ \\
\bottomrule
\end{tabular}
\end{center}

Summing and adding the atomic contributions, which by Theorem~\ref{thm:no-G0-atom}
occur only in $h_L$:

\begin{corollary}[Strict lift balance]
\label{cor:lift-balance}
For every $L \ge 7$,
\[
g_L = 2 g_{L-1} + h_{L-1}, \qquad h_L = h_{L-1} + |\Atom_L^{G_{\ne 0}}|.
\]
Adding, $|\Obs_L| = 2 |\Obs_{L-1}| + |\Atom_L^{G_{\ne 0}}|$.
\end{corollary}

\begin{theorem}\label{thm:series}
The density constant $c_W$ admits the convergent series representation
\[
c_W = \frac{3}{64} + \sum_{L \ge 7} \frac{|\Atom_L^{G_{\ne 0}}|}{2^L},
\]
with non-negative terms. In particular, $|\Atom_L^{G_{\ne 0}}|/2^L \to 0$ as
$L \to \infty$.
\end{theorem}

\begin{proof}
By Lemma~\ref{lem:atom-decomp} and Theorem~\ref{thm:no-G0-atom}, for $L \ge 7$:
\[
|\Obs_L| = 3 \cdot 2^{L-6} + \sum_{L_0 = 7}^{L} |\Atom_{L_0}^{G_{\ne 0}}| \cdot 2^{L - L_0},
\]
hence $\rho_L = 3/64 + \sum_{L_0 = 7}^{L} |\Atom_{L_0}^{G_{\ne 0}}|/2^{L_0}$.
By Theorem~\ref{thm:density-intro}, $\rho_L$ is non-decreasing and bounded
above by $1$, so its limit $c_W$ exists. The right-hand side is a partial sum
of a non-negative series; its limit is the announced series, which converges
to a finite value, forcing its terms to tend to $0$. \qedhere
\end{proof}

\subsection{The remaining open question}

Theorem~\ref{thm:series} reduces the question of $c_W$'s exact value to
controlling the growth of $|\Atom_L^{G_{\ne 0}}|$. A bound of the form
$|\Atom_L^{G_{\ne 0}}| \le C \cdot \alpha^L$ with $\alpha < 2$ would give a
geometrically convergent series with explicit error bounds. We do not
establish such a bound here; we return to the question in
Section~\ref{sec:open}.

% =============================================================================
\section{The \texorpdfstring{$T_+$}{T+} side}
\label{sec:tplus}
% =============================================================================

The classical Collatz map $\Tplus$ admits the same obstruction theory,
transferred from $\Tminus$ via the residue involution
$r \mapsto \overline{r} := (-r) \bmod 2^L$.

\begin{lemma}\label{lem:bijection}
For each odd $r$ and each level $L$, $\Xend^+(\overline{r}, L) = 1$
if and only if $\Xend(r, L) = 1$, where $\Xend^+$ is the $X$-invariant
defined as in Section~\ref{sec:setup} but with $\Tminus$ replaced by
$\Tplus$ throughout.
\end{lemma}

\begin{proof}
The basic $\Tplus$-step on $\overline{r}$ produces
$3\overline{r} + 1 = 3 \cdot 2^L - (3r - 1)$ in $\Z$.
The two summands have $2$-adic valuations $L$ and $\vtwo(3r - 1)$;
during the reduction the modulus has not yet been exhausted, so
$\vtwo(3r - 1) < L$ and the ultrametric identity yields
$v_+(\overline{r}) = \vtwo(3\overline{r} + 1) = \vtwo(3r - 1) = v_-(r)$.
The affine update therefore transfers without change, and the terminal
$X$-value is the same. \qedhere
\end{proof}

Write $\Obs_L^+$ for the $\Tplus$ analogue of $\Obs_L$.

\begin{theorem}\label{thm:transfer}
By Lemma~\ref{lem:bijection}, the map $r \mapsto \overline{r}$ restricts to a
bijection $\Obs_L \leftrightarrow \Obs_L^+$. Theorems~\ref{thm:lift-intro},~\ref{thm:density-intro},
and~\ref{thm:fc-intro} hold verbatim for $\Obs^+$, with the same density
constant $c_W$ and factor complexity function $p_W^+(n) = 2^n$.
\end{theorem}

% =============================================================================
\section{Open questions}
\label{sec:open}
% =============================================================================

\subsection{The exact value of \texorpdfstring{$c_W$}{cW}}

The principal open question is to determine $c_W$ exactly, or equivalently
to give a quantitative bound on the growth of $|\Atom_L^{G_{\ne 0}}|$. Three
approaches present themselves, each requiring substantial additional work.

\begin{sloppypar}
\textbf{(i) Substitution-theoretic.} Stérin~\cite{Sterin2020} realizes
$\Tplus$ as a four-state finite-state transducer between bases $2$ and $3$.
A related transducer-substitution system for the parallel reduction would,
combined with the classification of Pansiot~\cite{Pansiot1984} or
Devyatov~\cite{Devyatov2015}, determine the asymptotic growth class of
$\Atom_L^{G_{\ne 0}}$.
\end{sloppypar}

\textbf{(ii) Markov-operator-spectral.} Wirsching~\cite{Wirsching1998}
constructs a Markov operator $W_3$ on residue-density functions for the
$\Tplus$ predecessor problem. An analogous operator for the parallel
reduction would have an action on $(c, d)$-space; its spectrum would
determine the asymptotic growth of the shift-class populations
$|\Obs_L^{G_a}|$. A unified treatment --- for instance, expressing $c_W$
in terms of Wirsching's invariant density $\varphi$ --- appears not to be
in the literature.

\textbf{(iii) Generating-function-combinatorial.} A direct enumeration via
generating functions in $L$ is in principle accessible, but the absence of a
closed-form recursion --- the count depends on the joint distribution of
stopping times and shift indices --- makes the analysis delicate.

\subsection{Bridges to nontrivial cycles}

\begin{sloppypar}
The shift-zero class $\Obs_L^{G_0}$ encodes residues for which the two
parallel stopping times remain rigidly tied. Whether the existence or
non-existence of additional nontrivial cycles of $\Tminus$ is reflected in
the structure of $\Obs_L^{G_0}$ for large $L$ --- in particular, whether
one can extract a Diophantine obstruction from the strict lift balance of
Corollary~\ref{cor:lift-balance} --- is a natural and apparently unexplored
question. Classical lower bounds on the length of nontrivial $\Tplus$-cycles,
established by Eliahou~\cite{Eliahou1993} using continued-fraction methods
and extended numerically by Belaga--Mignotte~\cite{BelagaMignotte2006}, give
one indication of how rigidly the Diophantine geometry constrains cycle
structure; an analogous bound expressed in terms of $\Obs_L^{G_0}$ would be
of independent interest.
\end{sloppypar}

\subsection*{Acknowledgements}

The author declares no funding or competing interests.

% =============================================================================
\appendix
\section{Computational verification}
\label{sec:appendix}
% =============================================================================

The algorithmic $X$-invariant criterion of Section~\ref{sec:setup} is
implemented in Python in the supporting code that accompanies this paper,
collected in the directory \texttt{code/} alongside the manuscript source
and mirrored at
\url{https://github.com/yaccob/collatz/tree/main/manuscripts/t_minus/code}.
The directory contains the verification
scripts for each rigorous result (lift theorem, atom decomposition,
the no-shift-zero-atom theorem (Theorem~\ref{thm:no-G0-atom}), full factor
complexity, simultaneous stop) together with a
\texttt{README.md} mapping each script to the result it certifies. For
each level $L \le L_{\max}$, the relevant script enumerates all odd
residues $r$ with $\vtwo(r-1) \ge 1$, runs the parallel reduction, and
records the terminal data. The set $\Obs_L$ is read off as the preimage of
$\{(c, d) : 3d + c = 1\}$.

We have run this enumeration through $L = 16$, producing the following
counts:

\begin{center}
\begin{tabular}{rrrrr}
\toprule
$L$ & $|\Obs_L|$ & $|\Obs_L^{G_0}|$ & $|\Obs_L^{G_{\ne 0}}|$ & $|\Atom_L^{G_{\ne 0}}|$ \\
\midrule
$6$  & $3$    & $1$    & $2$   & $2$   \\
$7$  & $8$    & $4$    & $4$   & $2$   \\
$8$  & $19$   & $12$   & $7$   & $3$   \\
$9$  & $42$   & $31$   & $11$  & $4$   \\
$10$ & $91$   & $73$   & $18$  & $7$   \\
$11$ & $194$  & $164$  & $30$  & $12$  \\
$12$ & $409$  & $358$  & $51$  & $21$  \\
$13$ & $855$  & $767$  & $88$  & $37$  \\
$14$ & $1776$ & $1622$ & $154$ & $66$  \\
$15$ & $3671$ & $3398$ & $273$ & $119$ \\
$16$ & $7556$ & $7069$ & $487$ & $214$ \\
\bottomrule
\end{tabular}
\end{center}

The data satisfies the strict lift balance $|\Obs_L| = 2 |\Obs_{L-1}| + |\Atom_L^{G_{\ne 0}}|$
of Corollary~\ref{cor:lift-balance} at every level, and is consistent with
Theorems~\ref{thm:lift-intro},~\ref{thm:no-G0-atom}, and~\ref{thm:series}
(series convergence).

% =============================================================================
\bibliographystyle{amsalpha}
\bibliography{references}
% =============================================================================

\end{document}
